Put
\[
a_n=d\log(en)+\log(2n^2),
\]
and write \(W_i=(X_i,A_i)\) and \(\varepsilon_i=Y_i-r_\theta(W_i)\). By linear realizability,
\[
\mathbb E_P[\varepsilon_i\mid W_i]=0,
\]
while boundedness of \(Y_i\) and \(r_\theta\) gives \(|\varepsilon_i|\le1\). Conditional on \(W_1,\ldots,W_n\), the variables \(\varepsilon_1,\ldots,\varepsilon_n\) are independent. Let \(V_n\) be the subspace of \(\mathbb R^n\) consisting of vectors
\[
(f(W_1),\ldots,f(W_n)),\qquad f\in\operatorname{span}(\phi),
\]
and let \(\Pi_n\) be its orthogonal projector. The subspace \(V_n\) has dimension at most \(d\), and \(\Pi_n\) is measurable with respect to \(\sigma(W_1,\ldots,W_n)\). Indeed, it is the orthogonal projector generated by the sample feature matrix and hence is a Borel function of that matrix through the Moore--Penrose inverse.

Apply \(\text{lem:bounded-projection-concentration}\), conditionally on \(W_1,\ldots,W_n\), with failure probability \(\zeta=\delta_n/4\). After integrating the conditional probability, there is an event \(\mathcal E_{1,n}\) with
\[
P^{\otimes n}(\mathcal E_{1,n})\ge1-\frac{\delta_n}{4}
\]
on which
\[
\|\Pi_n\varepsilon\|_2^2
\le8\left\{d\log5+\log(8n^2)\right\}.
\]
Because \(r_\theta\in\mathcal F\), the empirical-risk-minimizing property of \(\widehat f\) gives
\[
\sum_{i=1}^n\{Y_i-\widehat f(W_i)\}^2
\le
\sum_{i=1}^n\{Y_i-r_\theta(W_i)\}^2.
\]
Writing \(h_i=\widehat f(W_i)-r_\theta(W_i)\), expansion yields
\[
\|h\|_2^2\le2\varepsilon^\top h.
\]
The vector \(h\) belongs to \(V_n\), so \(\varepsilon^\top h=(\Pi_n\varepsilon)^\top h\). Cauchy--Schwarz therefore gives
\[
\|h\|_2\le2\|\Pi_n\varepsilon\|_2
\]
and consequently, on \(\mathcal E_{1,n}\),
\[
\|\widehat f-r_\theta\|_n^2
\le
\frac{32}{n}\left\{d\log5+\log(8n^2)\right\}.
\]
For \(d\ge4\) and \(n\ge1\),
\[
d\log5+\log(8n^2)\le2a_n.
\]
Indeed, the difference between the right-hand side and the left-hand side equals
\[
d\{2+2\log n-\log5\}+2\log n-\log2,
\]
which is nonnegative at \(d=4,n=1\) and is increasing in both indices. Hence
\[
\|\widehat f-r_\theta\|_n^2\le64\frac{a_n}{n}.
\]
For \(c_0\ge64\), this proves \(r_\theta\in\mathcal Q_n\) on \(\mathcal E_{1,n}\).

We next transfer the empirical control to logging-population prediction error. The linear span containing \(\mathcal F\), equipped with the supremum norm on the finite set \(S\), has dimension at most \(d\). Since \(\mathcal F\subseteq[0,1]^S\), a volumetric argument gives
\[
N_{\mathcal F}(u)\le(1+2/u)^d,\qquad 0<u<1.
\]
To verify this bound, take a maximal \(u\)-separated subset of \(\mathcal F\). The supremum-norm balls of radius \(u/2\) around its points are disjoint and lie in the ball of radius \(1+u/2\) in the same at-most-\(d\)-dimensional linear space. Comparing volumes gives the displayed cardinality bound, while maximality makes the subset an \(u\)-cover.

Apply \(\text{lem:zhao-uniform-square-comparison}\) to \(\mathcal G=\mathcal F\), to the logging law \(\mu=\rho\times\pi_{\mathrm{ref}}\), with failure probability \(\delta_n/4\) and with
\[
u=\frac1{n+1}.
\]
There is an event \(\mathcal E_{2,n}\), of probability at least \(1-\delta_n/4\), on which simultaneously for every \(f_1,f_2\in\mathcal F\),
\[
\mathbb E_\mu(f_1-f_2)^2
\le
2\|f_1-f_2\|_n^2
+\frac{32}{3n}\log\{8n^2N_{\mathcal F}(u)\}
+\frac{10}{n+1}.
\]
The covering estimate gives
\[
N_{\mathcal F}(u)\le(2n+3)^d\le(5n)^d.
\]
Moreover,
\[
d\log(5n)\le2d\log(en)
\quad\text{and}\quad
\log(8n^2)\le3\log(2n^2).
\]
It follows that
\[
\log\{8n^2N_{\mathcal F}(u)\}\le3a_n.
\]
Since \(a_n\ge d\ge4\),
\[
\frac{32}{3n}\log\{8n^2N_{\mathcal F}(u)\}
+\frac{10}{n+1}
\le35\frac{a_n}{n}.
\]
Set \(\mathcal E_n=\mathcal E_{1,n}\cap\mathcal E_{2,n}\). The union bound gives
\[
P^{\otimes n}(\mathcal E_n)\ge1-\frac{\delta_n}{2}.
\]
On this event, for every \(f\in\mathcal Q_n\),
\[
\begin{aligned}
\|f-r_\theta\|_n^2
&\le2\|f-\widehat f\|_n^2+2\|\widehat f-r_\theta\|_n^2\\
&\le(2c_0+128)\frac{a_n}{n}.
\end{aligned}
\]
Taking \(f_1=f\) and \(f_2=r_\theta\) in the uniform comparison yields
\[
\begin{aligned}
\mathbb E_\mu(f-r_\theta)^2
&\le2(2c_0+128)\frac{a_n}{n}+35\frac{a_n}{n}\\
&=(4c_0+291)\frac{a_n}{n}.
\end{aligned}
\]
Taking the supremum over \(f\in\mathcal Q_n\) proves the general claim. With \(c_0=64\), the coefficient is \(547\), which is at most \(600\).

All optimization and event maps are measurable. The support \(S\) is finite, \(\mathcal F\) and \(\mathcal Q_n\) are compact finite-dimensional subsets of prediction space, and the stipulated lexicographic selection makes \(\widehat f\) a deterministic Borel function of the sample. The argument also covers \(n=1\): then \(\delta_1=1\), the asserted event probability is \(1/2\), and the comparison scale \(u=1/(n+1)=1/2\) remains in \((0,1)\).